\section{Literal certificates for the illustrated graphs}\label{sec:figures}
The following certificates complete the descriptive claims in the
source's Figures 2 and 4. They use the same three graphs as the positive
witness collection C1: \texttt{UntilOrder250\_50},
\texttt{UntilOrder250\_82} and \texttt{UntilOrder250\_300}, of orders
16, 20 and 40, respectively. All labels below are zero-based labels of
those exact graph6 records. An array $p$ lists $p(0),p(1),\ldots$;
line breaks within an array do not restart its indices. Complete arrays
and graph6 strings are also retained in the literal figure certificate.

\subsection{The Moebius--Kantor graph and its extension}
Define $G(8,3)$ on $u_i,v_i$, with $i$ taken modulo eight, by the edges
$u_iu_{i+1}$, $u_iv_i$ and $v_iv_{i+3}$. Label $u_i$ by $i$ and $v_i$
by $8+i$. The following full adjacency isomorphism identifies the
source's order-16 graph with this Moebius--Kantor graph:
\[
p=\bigl[\begin{gathered}
0,1,7,8,4,3,12,5,10,14,2,9,6,11,13,15
\end{gathered}\bigr].
\]
The permutations $u_i\mapsto u_{i+1}$, $v_i\mapsto v_{i+1}$
and $u_i\mapsto v_{3i}$, $v_i\mapsto u_{3i}$ preserve all three edge
types and act transitively together. Transporting them to the source
labels gives the two checked automorphisms
\[
a_1=\bigl[\begin{gathered}
1,10,0,11,7,4,14,12,13,15,5,8,2,6,9,3
\end{gathered}\bigr].
\]
\[
a_2=\bigl[\begin{gathered}
3,13,14,0,6,11,4,15,12,10,9,5,8,1,2,7
\end{gathered}\bigr].
\]
Their orbit of vertex zero contains all sixteen vertices. The following
injection identifies an induced copy of this graph in the source's
order-20 graph:
\[
\iota=\bigl[\begin{gathered}
1,8,17,16,2,19,18,11,3,4,9,10,14,12,15,13
\end{gathered}\bigr].
\]
The complementary set is $D=\{0,5,6,7\}$. It is independent;
every vertex outside $D$ has exactly one neighbour in $D$, and each
vertex of $D$ has four neighbours outside it. Checking all induced
adjacencies and these attachments verifies the whole add-four-vertices
construction in Figure 2. It also proves connectedness, since the
induced $G(8,3)$ is connected and each added vertex meets it.

Let $A$ be the adjacency matrix of this order-20 graph. Exact integer
multiplication gives
\[
 (A^6)_{vv}=\begin{cases}292,&v\in D,\\298,&v\notin D.\end{cases}
\]
This diagonal counts closed walks of length six and is invariant under
every automorphism. The two values therefore exclude vertex transitivity.
The literal cycle $(0,1,8,9,3)$ is a pentagon, so the graph is not bipartite.
To exclude edge transitivity, suppose a connected graph is edge-transitive.
The orbits of the two endpoints of one edge cover all vertices. If those
orbits coincide the graph is vertex-transitive. If they differ, every edge
joins the two orbits, giving a bipartition. Both alternatives fail here.
Thus the order-20 graph is neither vertex-transitive nor edge-transitive.

\subsection{The order-40 graph}
For the source's order-40 graph, the following three permutations are
checked on the entire adjacency relation:
\[
b_1=\bigl[\begin{gathered}
0,1,3,2,5,4,7,6,8,9,11,10,12,15,14,13,17,16,21,20,\\
19,18,23,22,33,31,30,32,29,28,26,25,27,24,34,38,39,37,35,36
\end{gathered}\bigr].
\]
\[
b_2=\bigl[\begin{gathered}
0,2,3,1,5,4,8,6,7,9,12,13,11,14,15,10,17,18,19,20,\\
21,16,24,25,26,27,28,29,32,30,33,22,31,23,35,36,37,39,34,38
\end{gathered}\bigr].
\]
\[
b_3=\bigl[\begin{gathered}
1,34,37,0,18,21,9,12,14,8,6,33,32,7,27,24,30,5,31,26,\\
4,25,10,16,20,15,38,36,11,17,13,35,19,39,22,29,3,23,2,28
\end{gathered}\bigr].
\]
Their generated orbit of vertex zero has forty vertices. This proves
the vertex-transitivity assertion of Figure 4. These are positive
automorphism and orbit certificates; no assertion that the displayed
generators generate the full automorphism group is needed.

All finite assertions in this section were independently checked from
the complete literal arrays, including the full closed-walk diagonal.
The separate companion arguments still supply the defining graph
parameters, uniqueness and extremality claims. The present certificates
do not infer those conclusions from symmetry or an embedding.
